
\documentclass{article}
\usepackage{colortbl}
\usepackage{makecell}
\usepackage{multirow}
\usepackage{supertabular}

\begin{document}

\newcounter{utterance}

\centering \large Interaction Transcript for game `wordle', experiment `medium\_frequency\_words\_no\_clue\_no\_critic', episode 10 with deepseek{-}v4{-}pro.
\vspace{24pt}

{ \footnotesize  \setcounter{utterance}{1}
\setlength{\tabcolsep}{0pt}
\begin{supertabular}{c@{$\;$}|p{.15\linewidth}@{}p{.15\linewidth}p{.15\linewidth}p{.15\linewidth}p{.15\linewidth}p{.15\linewidth}}
    \# & \multicolumn{2}{c}{Player} && \multicolumn{2}{c}{Game Master} \\
    \hline

    \theutterance \stepcounter{utterance}  
    & & & \multicolumn{4}{p{0.6\linewidth}}{
        \cellcolor[rgb]{0.9,0.9,0.9}{
            \makecell[{{p{\linewidth}}}]{
                \texttt{\tiny{[P1$\langle$GM]}}
                \texttt{Sa oled keelemaag, kellele meeldib sõnu ära arvata, kasutades etteantud reegleid.} \\
\\ 
\texttt{Tere tulemast Wordle'isse! Sul on kuus katset, et ära arvata sihtsõna — kehtiv viietäheline eestikeelne sõna väiketähtedest. Palun kasuta silte "selgitus:" ja "arvamus:", et anda iga pakkumise kohta lühike selgitus.} \\
\\ 
\texttt{Näiteks, kui sinu pakkumine on "aasta", peaks sinu vastus olema} \\
\texttt{selgitus: see on levinud viietäheline eestikeelne sõna ja alustan oma pakkumist selle sõnaga.} \\
\texttt{arvamus: aasta} \\
\\ 
\texttt{Pärast iga pakkumist kontrollitakse sinu vastust ning saad tagasisidet, mis näitab, millised tähed on õiged (roheline), millised tähed on õiged, kuid vales kohas (kollane), ja millised tähed on valed (punane). See tagasiside võib olla kasulik, et otsustada, milliseid tähti oma järgmises pakkumises kasutada või välistada.} \\
\\ 
\texttt{Näiteks tagasiside sõnale "aasta" võib olla:} \\
\texttt{arvamuse\_tagasiside: a<kollane> a<kollane> s<roheline> t<kollane> a<punane>} \\
\\ 
\texttt{Selgitus peaks sisaldama üksikasju selle kohta, kuidas arvamuse\_tagasiside'i kasutatakse uue pakkumiseni jõudmiseks.} \\
\\ 
\texttt{Ära genereeri mingit muud teksti.} \\
\\ 
\texttt{Alustame sinu esimese pakkumisega.} \\
            }
        }
    }
     \\ \\

    \theutterance \stepcounter{utterance}  
    & \multicolumn{4}{p{0.6\linewidth}}{
        \cellcolor[rgb]{0.561,0.737,0.561}{
            \makecell[{{p{\linewidth}}}]{
                \texttt{\tiny{[P1$\rangle$GM]}}
                \texttt{selgitus: Alustan sagedase ja täishäälikurikka eestikeelse sõnaga, et saada hea ülevaade võimalikest tähtedest.} \\
\texttt{arvamus: aedik} \\
            }
        }
    }
    & & \\ \\

    \theutterance \stepcounter{utterance}  
    & & & \multicolumn{2}{p{0.3\linewidth}}{
        \cellcolor[rgb]{0.9,0.9,0.9}{
            \makecell[{{p{\linewidth}}}]{
                \texttt{\tiny{[GM$|$GM]}}
                \texttt{attempts = 1} \\
\texttt{target\_word = meede} \\
\texttt{guess = aedik} \\
\texttt{guess\_feedback = a<red> e<green> d<yellow> i<red> k<red>} \\
            }
        }
    }
    & & \\ \\

    \theutterance \stepcounter{utterance}  
    & & & \multicolumn{4}{p{0.6\linewidth}}{
        \cellcolor[rgb]{0.9,0.9,0.9}{
            \makecell[{{p{\linewidth}}}]{
                \texttt{\tiny{[P1$\langle$GM]}}
                \texttt{arvamuse\_tagasiside: a<punane> e<roheline> d<kollane> i<punane> k<punane>} \\
\\ 
\texttt{Esitage oma vastus ainult selles vormingus.} \\
\texttt{selgitus: üksikasjad} \\
\texttt{arvamus: sõna} \\
            }
        }
    }
     \\ \\

    \theutterance \stepcounter{utterance}  
    & \multicolumn{4}{p{0.6\linewidth}}{
        \cellcolor[rgb]{0.561,0.737,0.561}{
            \makecell[{{p{\linewidth}}}]{
                \texttt{\tiny{[P1$\rangle$GM]}}
                \texttt{selgitus: Tean, et teine täht on E ja sõnas leidub D, mis pole kolmandal kohal. Välistan A, I ja K. Otsin uut sõna, mis sisaldaks E{-}d teisel kohal ja D{-}d mujal, kasutades levinumaid kaashäälikuid (R, L, S, T, N) ja täishäälikuid (U, O).} \\
\texttt{arvamus: redel} \\
            }
        }
    }
    & & \\ \\

    \theutterance \stepcounter{utterance}  
    & & & \multicolumn{2}{p{0.3\linewidth}}{
        \cellcolor[rgb]{0.9,0.9,0.9}{
            \makecell[{{p{\linewidth}}}]{
                \texttt{\tiny{[GM$|$GM]}}
                \texttt{attempts = 2} \\
\texttt{target\_word = meede} \\
\texttt{guess = redel} \\
\texttt{guess\_feedback = r<red> e<green> d<yellow> e<yellow> l<red>} \\
            }
        }
    }
    & & \\ \\

    \theutterance \stepcounter{utterance}  
    & & & \multicolumn{4}{p{0.6\linewidth}}{
        \cellcolor[rgb]{0.9,0.9,0.9}{
            \makecell[{{p{\linewidth}}}]{
                \texttt{\tiny{[P1$\langle$GM]}}
                \texttt{arvamuse\_tagasiside: r<punane> e<roheline> d<kollane> e<kollane> l<punane>} \\
\\ 
\texttt{Esitage oma vastus ainult selles vormingus.} \\
\texttt{selgitus: üksikasjad} \\
\texttt{arvamus: sõna} \\
            }
        }
    }
     \\ \\

    \theutterance \stepcounter{utterance}  
    & \multicolumn{4}{p{0.6\linewidth}}{
        \cellcolor[rgb]{0.561,0.737,0.561}{
            \makecell[{{p{\linewidth}}}]{
                \texttt{\tiny{[P1$\rangle$GM]}}
                \texttt{selgitus: Teisel kohal on E. Sõnas on D, mis ei ole kolmandal ega neljandal kohal. Uus ja teadaolevalt vales kohas täht on teine E. Välistan R, L. D on jäänud ainsaks võimalikuks kaashäälikuks esimesele või viimasele kohale. Proovin sõna, kus D on esimesel kohal ja E{-}d on teisel ning viimasel kohal.} \\
\texttt{arvamus: deeme} \\
            }
        }
    }
    & & \\ \\

    \theutterance \stepcounter{utterance}  
    & & & \multicolumn{2}{p{0.3\linewidth}}{
        \cellcolor[rgb]{0.9,0.9,0.9}{
            \makecell[{{p{\linewidth}}}]{
                \texttt{\tiny{[GM$|$GM]}}
                \texttt{Error: The guessed word is not a valid word for this game.} \\
            }
        }
    }
    & & \\ \\

    \theutterance \stepcounter{utterance}  
    & & & \multicolumn{4}{p{0.6\linewidth}}{
        \cellcolor[rgb]{0.9,0.9,0.9}{
            \makecell[{{p{\linewidth}}}]{
                \texttt{\tiny{[P1$\langle$GM]}}
                \texttt{Teie arvamus ei ole selle mängu jaoks kehtiv sõna. Palun proovige uuesti.} \\
\\ 
\texttt{Esitage oma vastus ainult selles vormingus.} \\
\texttt{selgitus: üksikasjad} \\
\texttt{arvamus: sõna} \\
            }
        }
    }
     \\ \\

    \theutterance \stepcounter{utterance}  
    & \multicolumn{4}{p{0.6\linewidth}}{
        \cellcolor[rgb]{0.561,0.737,0.561}{
            \makecell[{{p{\linewidth}}}]{
                \texttt{\tiny{[P1$\rangle$GM]}}
                \texttt{selgitus: Vabandust, "deeme" ei sobinud. Tean, et teine täht on E. Sõnas on D (mitte 3. ega 4. kohal) ja teine E (mitte 4. kohal). Välistatud on A, I, K, R, L. D asukoht on tõenäoliselt esimene. Kuna D on 1. kohal ja E on 2. kohal, proovin kombinatsiooni, kus viimane täht on E, et paigutada teine E õigesse kohta (kuna 4. koht on vale ja 2. koht on juba E). Otsin sõna kujul D E \_ \_ E.} \\
\texttt{arvamus: dedve} \\
            }
        }
    }
    & & \\ \\

    \theutterance \stepcounter{utterance}  
    & & & \multicolumn{2}{p{0.3\linewidth}}{
        \cellcolor[rgb]{0.9,0.9,0.9}{
            \makecell[{{p{\linewidth}}}]{
                \texttt{\tiny{[GM$|$GM]}}
                \texttt{Error: The guessed word is not a valid word for this game.} \\
            }
        }
    }
    & & \\ \\

    \theutterance \stepcounter{utterance}  
    & & & \multicolumn{4}{p{0.6\linewidth}}{
        \cellcolor[rgb]{0.9,0.9,0.9}{
            \makecell[{{p{\linewidth}}}]{
                \texttt{\tiny{[P1$\langle$GM]}}
                \texttt{Teie arvamus ei ole selle mängu jaoks kehtiv sõna. Palun proovige uuesti.} \\
\\ 
\texttt{Esitage oma vastus ainult selles vormingus.} \\
\texttt{selgitus: üksikasjad} \\
\texttt{arvamus: sõna} \\
            }
        }
    }
     \\ \\

    \theutterance \stepcounter{utterance}  
    & \multicolumn{4}{p{0.6\linewidth}}{
        \cellcolor[rgb]{0.561,0.737,0.561}{
            \makecell[{{p{\linewidth}}}]{
                \texttt{\tiny{[P1$\rangle$GM]}}
                \texttt{selgitus: Vabandust, "dedve" ei sobinud. Tean, et teine täht on E. Sõnas on D (mitte 3. ega 4. kohal) ja teine E (mitte 4. kohal). Välistatud on A, I, K, R, L. D asukoht on tõenäoliselt esimene. Kuna D on 1. kohal ja E on 2. kohal, proovin levinud sõna, kus on ainult üks E (teine E tagasisidest oli tegelikult juba paigas oleva E kollane peegeldus, kui see oleks olnud teine täht, aga kuna teine koht on E, siis üks E on juba roheline ja teine E kollane tähendab, et sõnas on kaks E{-}d). Kui sõnas on kaks E{-}d, siis teine E ei saa olla 4. kohal. Proovin sõna kujul D E \_ \_ \_, kus teine E on 3. või 5. kohal ja teine kaashäälik on levinud. Proovin "deebet".} \\
\texttt{arvamus: deebet} \\
            }
        }
    }
    & & \\ \\

    \theutterance \stepcounter{utterance}  
    & & & \multicolumn{2}{p{0.3\linewidth}}{
        \cellcolor[rgb]{0.9,0.9,0.9}{
            \makecell[{{p{\linewidth}}}]{
                \texttt{\tiny{[GM$|$GM]}}
                \texttt{Error: The length of the guessed word is not 5.} \\
            }
        }
    }
    & & \\ \\

    \theutterance \stepcounter{utterance}  
    & & & \multicolumn{2}{p{0.3\linewidth}}{
        \cellcolor[rgb]{0.9,0.9,0.9}{
            \makecell[{{p{\linewidth}}}]{
                \texttt{\tiny{[GM$|$GM]}}
                \texttt{game\_result = ABORT} \\
            }
        }
    }
    & & \\ \\

\end{supertabular}
}

\end{document}
